%!TEX root = main.tex

\section{Related Work}
\label{sec:related}

\paragraph{Instance and contrastive learning.}
Instance-level classification considers each image in a dataset as its own class~\cite{bojanowski2017unsupervised,dosovitskiy2016discriminative,wu2018unsupervised}.
Dosovitskiy~\etal~\cite{dosovitskiy2016discriminative} assign a class explicitly to each image and learn a linear classifier with as many classes as images in the dataset.
As this approach becomes quickly intractable, Wu~\etal~\cite{wu2018unsupervised} mitigate this issue by replacing the classifier with a memory bank that stores previously-computed representations.
They rely on noise contrastive estimation~\cite{gutmann2010noise} to compare instances, which is a special form of contrastive learning~\cite{hjelm2018learning,oord2018representation}.
He~\etal~\cite{he2019momentum} improve the training of contrastive methods by storing representations from a momentum encoder instead of the trained network.
More recently, Chen~\etal~\cite{chen2020simple} show that the memory bank can be entirely replaced with the elements from the same batch if the batch is large enough.
In contrast to this line of works, we avoid comparing every pair of images by mapping the image features to a set of trainable prototype vectors.

\paragraph{Clustering for deep representation learning.}
Our work is also related to clustering-based methods~\cite{asano2019self,bautista2016cliquecnn,caron2018deep,caron2019unsupervised,huang2019unsupervised,xie2016unsupervised,yang2016joint,zhuang2019local,gidaris2020learning,yan2020cluster}.
Caron~\etal~\cite{caron2018deep} show that $k$-means assignments can be used as pseudo-labels to learn visual representations.
This method scales to large uncurated dataset and can be used for pre-training of supervised networks~\cite{caron2019unsupervised}.
However, their formulation is not principled and recently, Asano~\etal~\cite{asano2019self} show how to cast the pseudo-label assignment problem as an instance of the optimal transport problem.
We consider a similar formulation to map representations to prototype vectors, but unlike~\cite{asano2019self} we keep the soft assignment produced by the Sinkhorn-Knopp algorithm~\cite{cuturi2013sinkhorn} instead of approximating it into a hard assignment.
Besides, unlike Caron~\etal~\cite{caron2018deep,caron2019unsupervised} and Asano~\etal~\cite{asano2019self}, we obtain online assignments which allows our method to scale gracefully to any dataset size.

\paragraph{Handcrafted pretext tasks.}
Many self-supervised methods manipulate the input data to extract a supervised signal in the form of a pretext task~\cite{doersch2015unsupervised,agrawal2015learning,jenni2018self,kim2018learning,larsson2016learning,mahendran2018cross,misra2016shuffle,pathak2017learning,pathak2016context,wang2015unsupervised,wang2017transitive,zhang2017split}.
We refer the reader to Jing~\etal~\cite{jing2019self} for an exhaustive and detailed review of this literature.
Of particular interest, Misra and van der Maaten~\cite{misra2019self} propose to encode the jigsaw puzzle task~\cite{noroozi2016unsupervised} as an invariant for contrastive learning.
Jigsaw tiles are non-overlapping crops with small resolution that cover only part (${\sim}20\%$) of the entire image area.
In contrast, our \texttt{multi-crop} strategy consists in simply sampling multiple random crops with two different sizes: a standard size and a smaller one.

